\documentclass[11pt]{article}
\usepackage{amsmath}
\usepackage{amssymb}
\usepackage{geometry}
\usepackage{booktabs}
\usepackage[hidelinks]{hyperref}
\geometry{a4paper, margin=1in}

% Every number in this manuscript is a macro read out of the frozen result JSON
% by tools/emit_frozen_numbers.py. No literal is typed, so a figure in the prose
% cannot survive a change in the artifact it came from.
% GENERATED FILE -- do not edit.
% Regenerate with: PYTHONPATH=src python3.12 tools/emit_frozen_numbers.py
% Source: results/official_fold/OFFICIAL_MULTI_METHOD_BOOTSTRAP_vs_P2RANK.json
%         data/cryptobench_apo/ALGEBRAIC_FIELD.json

\newcommand{\NUnits}{192}
\newcommand{\NBoot}{10000}
\newcommand{\BootSeed}{20260725}
\newcommand{\NRows}{2112}
\newcommand{\NMethods}{11}
\newcommand{\FisherSix}{0.760}
\newcommand{\FisherAll}{0.783}
\newcommand{\NCandidates}{9}
\newcommand{\NSelectFitUnits}{386}
\newcommand{\NSelectPickUnits}{384}
\newcommand{\SelectedArch}{parallel bank, integer-multiplicity fusion + multi-scale gate}
\newcommand{\SelectedArchAuc}{0.7446}
\newcommand{\NLinWires}{172}
\newcommand{\LinRidge}{1e-06}
\newcommand{\LinGateRadius}{18}
\newcommand{\LinGateWeight}{0.5}
\newcommand{\LinOperatingQ}{0.10}
\newcommand{\NTabWires}{645}
\newcommand{\NTabTables}{5152}
\newcommand{\TabWidth}{2}
\newcommand{\TabRounds}{16}
\newcommand{\NTabCells}{82432}
\newcommand{\NTabCellsEmpty}{825}
\newcommand{\TabCellsEmptyPct}{1.00}
\newcommand{\TabCellOccupancy}{14826}
\newcommand{\TabRidge}{0.03}
\newcommand{\TabCap}{32}
\newcommand{\NTabUsedFullFold}{4797}
\newcommand{\TabFanOut}{25682}
\newcommand{\NTabTrainRes}{234{,}838}
\newcommand{\TabGateRadius}{18}
\newcommand{\TabGateWeight}{1}
\newcommand{\TabOperatingQ}{0.09}
\newcommand{\TabCompileSeconds}{130}
\newcommand{\TabArtifactMB}{1.88}
\newcommand{\TabRuntime}{0.342}
\newcommand{\PTwoRuntime}{4.722}
\newcommand{\TabSpeedup}{13.8}
\newcommand{\TabPickAuc}{0.8045}
\newcommand{\NTabCandidates}{9}
\newcommand{\TabRidgeLow}{0.8045}
\newcommand{\TabRidgeLowValue}{0.03}
\newcommand{\TabRidgeMid}{0.8042}
\newcommand{\TabRidgeMidValue}{0.1}
\newcommand{\TabRidgeHigh}{0.8023}
\newcommand{\TabRidgeHighValue}{0.3}
\newcommand{\CaseLocateF}{0.5}
\newcommand{\NLocBoth}{25}
\newcommand{\NLocOurs}{24}
\newcommand{\NLocPTwo}{13}
\newcommand{\NLocNeither}{130}
\newcommand{\CaseBoth}{1zm0\_A}
\newcommand{\CaseBothFOurs}{0.80}
\newcommand{\CaseBothFPTwo}{0.80}
\newcommand{\CaseBothOffOurs}{1}
\newcommand{\CaseBothOffPTwo}{2}
\newcommand{\CaseOurs}{2d05\_A}
\newcommand{\CaseOursFOurs}{0.89}
\newcommand{\CaseOursFPTwo}{0.29}
\newcommand{\CaseOursOffOurs}{2}
\newcommand{\CaseOursOffPTwo}{7}
\newcommand{\CaseTheirs}{6bty\_B}
\newcommand{\CaseTheirsFOurs}{0.10}
\newcommand{\CaseTheirsFPTwo}{0.57}
\newcommand{\CaseTheirsOffOurs}{11}
\newcommand{\CaseTheirsOffPTwo}{2}
\newcommand{\CaseNeither}{1bzj\_A}
\newcommand{\CaseNeitherFOurs}{0.05}
\newcommand{\CaseNeitherFPTwo}{0.13}
\newcommand{\CaseNeitherOffOurs}{12}
\newcommand{\CaseNeitherOffPTwo}{14}
\newcommand{\BurialBoth}{+1.31}
\newcommand{\BurialOurs}{+1.50}
\newcommand{\BurialTheirs}{+0.73}
\newcommand{\BurialNeither}{+0.35}
\newcommand{\BurialCallsNeitherOurs}{+0.86}
\newcommand{\BurialCallsNeitherPTwo}{+1.37}
\newcommand{\FigCapOfficialFoldMetrics}{Per-residue detection on the official CryptoBench apo test fold (192 single-chain units, MMseqs2 10\% cluster-disjoint). Bars are means over structures, ordered by ROC-AUC; whiskers are 95\% bootstrap intervals (10,000 resamples, seed 20260725). The table field is drawn in orange and P2Rank 2.5.1, the baseline, in blue. MCC is undefined for a detector that calls nothing positive; those entries are marked rather than left blank.}
\newcommand{\FigCapPairedVsPTwoRank}{Paired differences against P2Rank, on the structures where both are defined: one row per metric, showing the point estimate, the 95\% bootstrap interval and the number of shared structures. The shaded band is $\pm$1 standard error of a fold mean (0.0161); 12 of our architectures have been scored on this fold, so a margin inside that band is not an ordering. A row is called resolved only where the interval clears zero and the verdict survives every resampling seed, which is why three of the four are not.}
\newcommand{\FigCapFunctionalChoice}{One set of 192 paired differences, and what follows from how it is summarised. Left, the differences themselves: the field is ahead on 115 structures and behind on 77, a win rate +0.0990 above one half (p=0.004), and the shaded tails are what a 20\% trim discards --- losses averaging -0.2319 against wins of +0.1757. That asymmetry is why the mean sits at +0.0058 while the middle 60\% sits at +0.0281. Right, five location summaries of those same numbers with 95\% bootstrap intervals. The 20\% trimmed mean was chosen on the training partition and committed in 39e70d4 before this fold was read for it; the others are shown because reporting only the chosen one would hide whether the choice landed on the flattering summary. The mean is unresolved and stays in the paper: it is the functional that speaks to the discarded tail, where the field's failures are worse than P2Rank's.}
\newcommand{\FigCapMatchedOperatingPoint}{The F1 margin over P2Rank, before and after holding the operating point in common, on all 192 held-out structures. Left, mean per-structure F1 as a function of the fraction q of each chain called positive, for both methods. The dashed level is P2Rank's own pocket assignment, which is not a point on the curve because it is not a fraction of the chain; the published comparison is that level against our curve at q=0.09. Running the grid search that chose our q on P2Rank's training scores returns the same q=0.09, so the two matched rules a reader might ask for coincide. Right, the paired difference under each convention with 95\% bootstrap intervals: +0.0315 at the shipped operating points, +0.0140 [-0.0085, +0.0364] once both are binarised the same way, and +0.0108 against P2Rank at the best q this fold admits --- an upper bound it could not have known. Our own F1 is identical in the first two columns: our shipped call already is the matched rule at this q. The rules and the sentence to be written under each outcome were committed in c5911a3 before the fold was read for them.}
\newcommand{\FigCapMatchedFullMetrics}{Where the reported margin came from, on all 192 held-out structures. Left, the two detectors in the precision--recall plane under three conventions, each pair joined by a line. Run as delivered, P2Rank calls 1.65 times as many residues as our rule does (8564 against 5202), which buys recall and costs precision: our precision is higher by +0.0590 and our recall lower by -0.0662, both intervals excluding zero. At a common budget both call 5202 residues and the trade is gone. Right, every paired difference with a 95\% bootstrap interval over structures. The F1 and MCC differences remain positive and contain zero under every matched convention. The matched recall difference, +0.0332 [+0.0027, +0.0646], excludes zero at p=0.033, which does not survive a Bonferroni threshold of 0.0125 over the four metrics of one convention and carried no preregistered decision rule; it is an observation about this fold and is drawn as one. Resampling PDB entries or MMseqs2 clusters instead of chains changes no interval here by more than a factor of 1.006.}
\newcommand{\FigCapCaseStudies}{Four evaluation units, projected onto each chain's two longest principal axes: grey is the chain, black rings are the labelled cryptic residues, filled points are the residues each method called positive. The cases are chosen by rule and not by eye -- the clearest joint success, the largest margin each way, and the largest labelled pocket among the units neither method locates -- where locating means per-residue F1 at least 0.5. That last case is the commonest outcome: of 192 units, 25 are located by both methods, 24 by the table field only, 13 by P2Rank only, and 130 by neither.}
\newcommand{\NSplitsCv}{4}
\newcommand{\NFirstCv}{4}
\newcommand{\WorstRankCv}{1}
\newcommand{\MeanAucCv}{0.7392}
\newcommand{\SdAucCv}{0.0137}
\newcommand{\MeanMarginCv}{+0.0043}
\newcommand{\WorstMarginCv}{+0.0033}
\newcommand{\NSplitsRh}{25}
\newcommand{\NFirstRh}{25}
\newcommand{\WorstRankRh}{1}
\newcommand{\MeanAucRh}{0.7382}
\newcommand{\SdAucRh}{0.0099}
\newcommand{\MeanMarginRh}{+0.0055}
\newcommand{\WorstMarginRh}{+0.0032}
\newcommand{\RunnerUpArch}{parallel bank, uniform fusion + multi-scale gate}
\newcommand{\PreRegNPick}{384}
\newcommand{\PreRegFieldPick}{0.8045}
\newcommand{\PreRegPTwoPick}{0.7700}
\newcommand{\PreRegMeanPowerQuarter}{0.03}
\newcommand{\PreRegTrimPowerQuarter}{0.56}
\newcommand{\PreRegStratPowerQuarter}{0.03}
\newcommand{\PreRegChosen}{trimmed20}
\newcommand{\PreRegForecast}{39}
\newcommand{\PreRegCommit}{\texttt{39e70d43547a}}
\newcommand{\PreRegReadIndex}{5}
\newcommand{\PreRegTestDelta}{+0.0281}
\newcommand{\PreRegTestCI}{[+0.0117, +0.0427]}
\newcommand{\PreRegTestP}{0.002}
\newcommand{\PreRegTestMedian}{+0.0412}
\newcommand{\PreRegTestMedianP}{0.006}
\newcommand{\PreRegTestWinRate}{+0.0990}
\newcommand{\PreRegTestWinRateP}{0.004}
\newcommand{\PreRegNAhead}{115}
\newcommand{\PreRegNBehind}{77}
\newcommand{\PreRegNTrimmed}{38}
\newcommand{\PreRegTailLoss}{-0.2319}
\newcommand{\PreRegTailWin}{+0.1757}
\newcommand{\PTwoTrainAuc}{0.7714}
\newcommand{\PTwoTrainUnits}{770}
\newcommand{\MatchPTwoQ}{0.09}
\newcommand{\MatchPTwoTrainNative}{0.2956}
\newcommand{\MatchPTwoTrainTuned}{0.3203}
\newcommand{\MatchTrainGain}{+0.0248}
\newcommand{\MatchForecast}{+0.0068}
\newcommand{\MatchCommit}{\texttt{c5911a3bea5f}}
\newcommand{\MatchReadIndex}{6}
\newcommand{\MatchNUnits}{192}
\newcommand{\MatchOurQ}{0.09}
\newcommand{\MatchNativeOurs}{0.3334}
\newcommand{\MatchNativePTwo}{0.3019}
\newcommand{\MatchNativeDelta}{+0.0315}
\newcommand{\MatchAOurs}{0.3334}
\newcommand{\MatchAPTwo}{0.3194}
\newcommand{\MatchADelta}{+0.0140}
\newcommand{\MatchACI}{[-0.0085, +0.0364]}
\newcommand{\MatchAP}{0.224}
\newcommand{\MatchATrim}{+0.0192}
\newcommand{\MatchATrimCI}{[-0.0108, +0.0471]}
\newcommand{\MatchBOurs}{0.3334}
\newcommand{\MatchBPTwo}{0.3194}
\newcommand{\MatchBDelta}{+0.0140}
\newcommand{\MatchBCI}{[-0.0085, +0.0364]}
\newcommand{\MatchBP}{0.224}
\newcommand{\MatchBTrim}{+0.0192}
\newcommand{\MatchBTrimCI}{[-0.0108, +0.0471]}
\newcommand{\MatchOracleQ}{0.10}
\newcommand{\MatchOraclePTwo}{0.3226}
\newcommand{\MatchOracleDelta}{+0.0108}
\newcommand{\MatchOracleCI}{[-0.0113, +0.0326]}
\newcommand{\MatchOurOracleQ}{0.11}
\newcommand{\MatchOurOracleGap}{0.0022}
\newcommand{\FullDeployedCalledOurs}{5202}
\newcommand{\FullDeployedCalledPTwo}{8564}
\newcommand{\FullDeployedPrecOurs}{0.3177}
\newcommand{\FullDeployedPrecPTwo}{0.2591}
\newcommand{\FullDeployedPrecDelta}{+0.0590}
\newcommand{\FullDeployedPrecCI}{[+0.0363, +0.0813]}
\newcommand{\FullDeployedPrecP}{0.000}
\newcommand{\FullDeployedPrecN}{186}
\newcommand{\FullDeployedRecOurs}{0.4343}
\newcommand{\FullDeployedRecPTwo}{0.5005}
\newcommand{\FullDeployedRecDelta}{-0.0662}
\newcommand{\FullDeployedRecCI}{[-0.0999, -0.0318]}
\newcommand{\FullDeployedRecP}{0.000}
\newcommand{\FullDeployedRecN}{192}
\newcommand{\FullDeployedFOneOurs}{0.3334}
\newcommand{\FullDeployedFOnePTwo}{0.3019}
\newcommand{\FullDeployedFOneDelta}{+0.0315}
\newcommand{\FullDeployedFOneCI}{[+0.0089, +0.0537]}
\newcommand{\FullDeployedFOneP}{0.006}
\newcommand{\FullDeployedFOneN}{192}
\newcommand{\FullDeployedMCCOurs}{0.3040}
\newcommand{\FullDeployedMCCPTwo}{0.2848}
\newcommand{\FullDeployedMCCDelta}{+0.0247}
\newcommand{\FullDeployedMCCCI}{[+0.0004, +0.0494]}
\newcommand{\FullDeployedMCCP}{0.046}
\newcommand{\FullDeployedMCCN}{186}
\newcommand{\FullBudgetQOurs}{0.09}
\newcommand{\FullBudgetQPTwo}{0.09}
\newcommand{\FullBudgetCalledOurs}{5202}
\newcommand{\FullBudgetCalledPTwo}{5202}
\newcommand{\FullBudgetPrecOurs}{0.3177}
\newcommand{\FullBudgetPrecPTwo}{0.3114}
\newcommand{\FullBudgetPrecDelta}{+0.0063}
\newcommand{\FullBudgetPrecCI}{[-0.0162, +0.0280]}
\newcommand{\FullBudgetPrecP}{0.568}
\newcommand{\FullBudgetPrecN}{192}
\newcommand{\FullBudgetRecOurs}{0.4343}
\newcommand{\FullBudgetRecPTwo}{0.4010}
\newcommand{\FullBudgetRecDelta}{+0.0332}
\newcommand{\FullBudgetRecCI}{[+0.0027, +0.0646]}
\newcommand{\FullBudgetRecP}{0.033}
\newcommand{\FullBudgetRecN}{192}
\newcommand{\FullBudgetFOneOurs}{0.3334}
\newcommand{\FullBudgetFOnePTwo}{0.3194}
\newcommand{\FullBudgetFOneDelta}{+0.0140}
\newcommand{\FullBudgetFOneCI}{[-0.0085, +0.0364]}
\newcommand{\FullBudgetFOneP}{0.224}
\newcommand{\FullBudgetFOneN}{192}
\newcommand{\FullBudgetMCCOurs}{0.3040}
\newcommand{\FullBudgetMCCPTwo}{0.2871}
\newcommand{\FullBudgetMCCDelta}{+0.0169}
\newcommand{\FullBudgetMCCCI}{[-0.0092, +0.0424]}
\newcommand{\FullBudgetMCCP}{0.203}
\newcommand{\FullBudgetMCCN}{192}
\newcommand{\FullTunedFQOurs}{0.09}
\newcommand{\FullTunedFQPTwo}{0.09}
\newcommand{\FullTunedFCalledOurs}{5202}
\newcommand{\FullTunedFCalledPTwo}{5202}
\newcommand{\FullTunedFPrecOurs}{0.3177}
\newcommand{\FullTunedFPrecPTwo}{0.3114}
\newcommand{\FullTunedFPrecDelta}{+0.0063}
\newcommand{\FullTunedFPrecCI}{[-0.0162, +0.0280]}
\newcommand{\FullTunedFPrecP}{0.568}
\newcommand{\FullTunedFPrecN}{192}
\newcommand{\FullTunedFRecOurs}{0.4343}
\newcommand{\FullTunedFRecPTwo}{0.4010}
\newcommand{\FullTunedFRecDelta}{+0.0332}
\newcommand{\FullTunedFRecCI}{[+0.0027, +0.0646]}
\newcommand{\FullTunedFRecP}{0.033}
\newcommand{\FullTunedFRecN}{192}
\newcommand{\FullTunedFFOneOurs}{0.3334}
\newcommand{\FullTunedFFOnePTwo}{0.3194}
\newcommand{\FullTunedFFOneDelta}{+0.0140}
\newcommand{\FullTunedFFOneCI}{[-0.0085, +0.0364]}
\newcommand{\FullTunedFFOneP}{0.224}
\newcommand{\FullTunedFFOneN}{192}
\newcommand{\FullTunedFMCCOurs}{0.3040}
\newcommand{\FullTunedFMCCPTwo}{0.2871}
\newcommand{\FullTunedFMCCDelta}{+0.0169}
\newcommand{\FullTunedFMCCCI}{[-0.0092, +0.0424]}
\newcommand{\FullTunedFMCCP}{0.203}
\newcommand{\FullTunedFMCCN}{192}
\newcommand{\FullTunedMQOurs}{0.11}
\newcommand{\FullTunedMQPTwo}{0.09}
\newcommand{\FullTunedMCalledOurs}{6365}
\newcommand{\FullTunedMCalledPTwo}{5202}
\newcommand{\FullTunedMPrecOurs}{0.2955}
\newcommand{\FullTunedMPrecPTwo}{0.3114}
\newcommand{\FullTunedMPrecDelta}{-0.0159}
\newcommand{\FullTunedMPrecCI}{[-0.0380, +0.0054]}
\newcommand{\FullTunedMPrecP}{0.140}
\newcommand{\FullTunedMPrecN}{192}
\newcommand{\FullTunedMRecOurs}{0.4878}
\newcommand{\FullTunedMRecPTwo}{0.4010}
\newcommand{\FullTunedMRecDelta}{+0.0867}
\newcommand{\FullTunedMRecCI}{[+0.0528, +0.1212]}
\newcommand{\FullTunedMRecP}{0.000}
\newcommand{\FullTunedMRecN}{192}
\newcommand{\FullTunedMFOneOurs}{0.3356}
\newcommand{\FullTunedMFOnePTwo}{0.3194}
\newcommand{\FullTunedMFOneDelta}{+0.0162}
\newcommand{\FullTunedMFOneCI}{[-0.0069, +0.0391]}
\newcommand{\FullTunedMFOneP}{0.172}
\newcommand{\FullTunedMFOneN}{192}
\newcommand{\FullTunedMMCCOurs}{0.3075}
\newcommand{\FullTunedMMCCPTwo}{0.2871}
\newcommand{\FullTunedMMCCDelta}{+0.0204}
\newcommand{\FullTunedMMCCCI}{[-0.0065, +0.0472]}
\newcommand{\FullTunedMMCCP}{0.136}
\newcommand{\FullTunedMMCCN}{192}
\newcommand{\FullCallRatio}{1.65}
\newcommand{\FullGroupsChain}{192}
\newcommand{\FullGroupsPdb}{190}
\newcommand{\FullGroupsClu}{190}
\newcommand{\FullWidestClusterRatio}{1.006}
\newcommand{\FullIntervals}{48}
\newcommand{\FullBonfAlpha}{0.0125}
\newcommand{\FullReadIndex}{7}
\newcommand{\TrainGainF}{+0.0248}
\newcommand{\TrainGainM}{+0.0120}
\newcommand{\TrainQOurMCC}{0.11}
\newcommand{\TrainQPTwoMCC}{0.09}
\newcommand{\TrainQUnchangedOutOfSample}{yes}
\newcommand{\FullForecastFOne}{+0.0068}
\newcommand{\FullForecastMCC}{+0.0127}
\newcommand{\GenOperator}{190}
\newcommand{\GenChain}{77}
\newcommand{\GenTotal}{267}
\newcommand{\GenCeilBase}{0.7595}
\newcommand{\GenCeilFull}{0.7676}
\newcommand{\GenCeilLift}{+0.0081}
\newcommand{\GenCeilWorst}{+0.0039}
\newcommand{\GenNSplits}{12}
\newcommand{\GenNPositive}{12}
\newcommand{\GenFamTrace}{+0.0043}
\newcommand{\GenFamGyration}{+0.0036}
\newcommand{\GenFamDiagonal}{+0.0008}
\newcommand{\GenFamHinge}{+0.0008}
\newcommand{\GenFamShape}{+0.0027}
\newcommand{\GenFamLag}{+0.0021}
\newcommand{\GenFamValuation}{+0.0023}
\newcommand{\GenBestFamily}{spectral trace functional}
\newcommand{\GenWorstFamily}{soft-mode hinge}
\newcommand{\GenSecondWorstFamily}{diagonal functional at the centre}
\newcommand{\GenWires}{1713}
\newcommand{\GenWiresBase}{645}
\newcommand{\GenWireAuc}{0.8035}
\newcommand{\GenWireControlAuc}{0.8045}
\newcommand{\GenWireDelta}{-0.0010}
\newcommand{\GenWireRidge}{0.3}
\newcommand{\GenWireControlRidge}{0.03}
\newcommand{\GenWireEmpty}{4.17}
\newcommand{\GenWireControlEmpty}{1.25}
\newcommand{\SensLevelsThree}{0.7995}
\newcommand{\SensLevelsThreePooled}{0.7950}
\newcommand{\SensLevelsThreeEmpty}{0.87}
\newcommand{\SensLevelsFour}{0.8045}
\newcommand{\SensLevelsFourPooled}{0.7905}
\newcommand{\SensLevelsFourEmpty}{1.25}
\newcommand{\SensLevelsFive}{0.7996}
\newcommand{\SensLevelsFivePooled}{0.7961}
\newcommand{\SensLevelsFiveEmpty}{2.28}
\newcommand{\SensCapSixteen}{0.8018}
\newcommand{\SensCapThirtyTwo}{0.8045}
\newcommand{\SensCapSixtyFour}{0.8010}
\newcommand{\SensLevelsThreeDelta}{0.0050}
\newcommand{\SensLevelsFiveDelta}{0.0049}
\newcommand{\NTabUsedFitHalf}{4853}
\newcommand{\TabFanOutFitHalf}{29075}
\newcommand{\NTabCellsEmptyFitHalf}{1028}
\newcommand{\SensRange}{0.0140}
\newcommand{\SensSettings}{8}
\newcommand{\SensWorst}{0.7905}
\newcommand{\SensBest}{0.8045}
\newcommand{\NTrainPositives}{13,524}
\newcommand{\QuoDenseBound}{6.86}
\newcommand{\QuoDenseWidth}{6}
\newcommand{\QuoSymWidth}{41}
\newcommand{\QuoDenseCells}{4,096}
\newcommand{\QuoSymCells}{84}
\newcommand{\QuoSymCellsEight}{1,716}
\newcommand{\QuoSymCellsAll}{8,436}
\newcommand{\QuoNSplits}{14}
\newcommand{\QuoNCandidates}{10}
\newcommand{\QuoTrainGain}{+0.0087}
\newcommand{\QuoTrainWorst}{+0.0024}
\newcommand{\QuoNBeating}{14}
\newcommand{\QuoOneSplitGain}{+0.0132}
\newcommand{\QuoWorstArchDelta}{-0.1390}
\newcommand{\QuoReadIndex}{4}
\newcommand{\QuoTestAuc}{0.7647}
\newcommand{\QuoNTables}{18}
\newcommand{\QuoControlRepro}{0.7667}
\newcommand{\QuoDAf}{-0.0020}
\newcommand{\QuoDAfCI}{[-0.0103, +0.0065]}
\newcommand{\QuoDAfP}{0.6336}
\newcommand{\QuoDPTwo}{-0.0286}
\newcommand{\QuoDPTwoCI}{[-0.0486, -0.0095]}
\newcommand{\QuoDPTwoP}{0.0034}
\newcommand{\GapTablesSmall}{0.7446}
\newcommand{\GapLinearSmall}{0.7702}
\newcommand{\GapLinearWide}{0.7880}
\newcommand{\GapTablesWide}{0.7413}
\newcommand{\GapWideWires}{172}
\newcommand{\GapExtraWires}{137}
\newcommand{\GapWideTables}{29}
\newcommand{\GapReadout}{+0.0257}
\newcommand{\GapInput}{+0.0177}
\newcommand{\GapTotal}{+0.0434}
\newcommand{\GapReadoutShare}{59}
\newcommand{\GapInputShare}{41}
\newcommand{\GapFanoutPrice}{+0.0152}
\newcommand{\GapFanoutSolved}{0.7598}
\newcommand{\GapFanoutQuo}{-0.0015}
\newcommand{\GapScaleSmallPos}{878}
\newcommand{\GapScaleSmallGain}{+0.0094}
\newcommand{\GapScaleBigPos}{6620}
\newcommand{\GapScaleBigGain}{+0.0132}
\newcommand{\GapMarginalLevels}{64}
\newcommand{\GapMarginalAuc}{0.7001}
\newcommand{\GapMarginalTables}{35}
\newcommand{\GapUnseenPct}{0.15}
\newcommand{\GapHardCut}{0.55}
\newcommand{\GapHardTrainN}{70}
\newcommand{\GapHardTrainGain}{+0.0546}
\newcommand{\GapHardTestN}{28}
\newcommand{\GapHardTestGain}{+0.0277}
\newcommand{\GapEasyTrainGain}{-0.0023}
\newcommand{\GapEasyTestGain}{-0.0108}
\newcommand{\GapReweighted}{+0.0000}
\newcommand{\GapPickUnits}{384}
\newcommand{\NArchOnFold}{12}
\newcommand{\NOurDetectors}{9}
\newcommand{\NStandaloneProbes}{8}
\newcommand{\FoldSE}{0.0161}
\newcommand{\LeadInSE}{0.36}
\newcommand{\ReadAucI}{0.7804}
\newcommand{\ReadDeltaI}{-0.0131}
\newcommand{\ReadAucII}{0.7952}
\newcommand{\ReadDeltaII}{+0.0017}
\newcommand{\NTestReads}{3}
\newcommand{\NSeedProbes}{25}
\newcommand{\TabAucSeedSig}{0}
\newcommand{\TabPrSeedSig}{0}
\newcommand{\TabMccSeedSig}{6}
\newcommand{\TabFOneSeedSig}{25}
\newcommand{\PickCount}{0.7631}
\newcommand{\PickLinDigits}{0.7873}
\newcommand{\PickLinOnehot}{0.7844}
\newcommand{\PickLinTables}{0.7633}
\newcommand{\PickPairBest}{0.7846}
\newcommand{\PickPairBestName}{\texttt{pair\_30+lin}}
\newcommand{\NTrainUnits}{770}
\newcommand{\NTrainResidues}{234838}
\newcommand{\TrainRate}{0.0576}
\newcommand{\NAlgFeatures}{35}
\newcommand{\NAlgTables}{6}
\newcommand{\AlgOperatingQ}{0.09}
\newcommand{\LedgerTestClusters}{190}
\newcommand{\LedgerTestUnits}{192}
\newcommand{\LedgerTrainClusters}{766}
\newcommand{\LedgerTrainUnits}{770}

\newcommand{\TabAuc}{0.799}
\newcommand{\TabAucCI}{[0.766, 0.830]}
\newcommand{\AlgAuc}{0.767}
\newcommand{\AlgAucCI}{[0.736, 0.796]}
\newcommand{\AlgLinAuc}{0.795}
\newcommand{\AlgLinAucCI}{[0.765, 0.824]}
\newcommand{\PtwoRAuc}{0.793}
\newcommand{\PtwoRAucCI}{[0.772, 0.814]}
\newcommand{\GeoAuc}{0.665}
\newcommand{\GeoAucCI}{[0.644, 0.685]}
\newcommand{\QlutAuc}{0.760}
\newcommand{\QlutAucCI}{[0.733, 0.785]}
\newcommand{\QlutSeqAuc}{0.767}
\newcommand{\QlutSeqAucCI}{[0.740, 0.792]}
\newcommand{\SstarAuc}{0.655}
\newcommand{\SstarAucCI}{[0.634, 0.676]}
\newcommand{\FstarAuc}{0.599}
\newcommand{\FstarAucCI}{[0.579, 0.618]}
\newcommand{\UltraAuc}{0.653}
\newcommand{\UltraAucCI}{[0.632, 0.673]}
\newcommand{\RandAuc}{0.500}
\newcommand{\RandAucCI}{[0.500, 0.500]}
\newcommand{\TabAucD}{+0.0058}
\newcommand{\TabAucDCI}{[-0.0168, +0.0276]}
\newcommand{\TabAucP}{0.598}
\newcommand{\TabAucNPaired}{192}
\newcommand{\TabAucMatched}{0.799}
\newcommand{\TabAucMatchedBase}{0.793}
\newcommand{\TabAucVerdict}{indistinguishable}
\newcommand{\TabAucNS}{~ns}
\newcommand{\AlgAucD}{-0.0266}
\newcommand{\AlgAucDCI}{[-0.0487, -0.0055]}
\newcommand{\AlgAucP}{0.012}
\newcommand{\AlgAucNPaired}{192}
\newcommand{\AlgAucMatched}{0.767}
\newcommand{\AlgAucMatchedBase}{0.793}
\newcommand{\AlgAucVerdict}{separated}
\newcommand{\AlgAucNS}{~\textbf{sig}}
\newcommand{\AlgLinAucD}{+0.0020}
\newcommand{\AlgLinAucDCI}{[-0.0190, +0.0221]}
\newcommand{\AlgLinAucP}{0.841}
\newcommand{\AlgLinAucNPaired}{192}
\newcommand{\AlgLinAucMatched}{0.795}
\newcommand{\AlgLinAucMatchedBase}{0.793}
\newcommand{\AlgLinAucVerdict}{indistinguishable}
\newcommand{\AlgLinAucNS}{~ns}
\newcommand{\GeoAucD}{-0.1287}
\newcommand{\GeoAucDCI}{[-0.1461, -0.1111]}
\newcommand{\GeoAucP}{<0.001}
\newcommand{\GeoAucNPaired}{192}
\newcommand{\GeoAucMatched}{0.665}
\newcommand{\GeoAucMatchedBase}{0.793}
\newcommand{\GeoAucVerdict}{separated}
\newcommand{\GeoAucNS}{~\textbf{sig}}
\newcommand{\QlutAucD}{-0.0334}
\newcommand{\QlutAucDCI}{[-0.0517, -0.0160]}
\newcommand{\QlutAucP}{<0.001}
\newcommand{\QlutAucNPaired}{192}
\newcommand{\QlutAucMatched}{0.760}
\newcommand{\QlutAucMatchedBase}{0.793}
\newcommand{\QlutAucVerdict}{separated}
\newcommand{\QlutAucNS}{~\textbf{sig}}
\newcommand{\QlutSeqAucD}{-0.0268}
\newcommand{\QlutSeqAucDCI}{[-0.0442, -0.0099]}
\newcommand{\QlutSeqAucP}{0.002}
\newcommand{\QlutSeqAucNPaired}{192}
\newcommand{\QlutSeqAucMatched}{0.767}
\newcommand{\QlutSeqAucMatchedBase}{0.793}
\newcommand{\QlutSeqAucVerdict}{separated}
\newcommand{\QlutSeqAucNS}{~\textbf{sig}}
\newcommand{\SstarAucD}{-0.1383}
\newcommand{\SstarAucDCI}{[-0.1570, -0.1194]}
\newcommand{\SstarAucP}{<0.001}
\newcommand{\SstarAucNPaired}{192}
\newcommand{\SstarAucMatched}{0.655}
\newcommand{\SstarAucMatchedBase}{0.793}
\newcommand{\SstarAucVerdict}{separated}
\newcommand{\SstarAucNS}{~\textbf{sig}}
\newcommand{\FstarAucD}{-0.1946}
\newcommand{\FstarAucDCI}{[-0.2163, -0.1722]}
\newcommand{\FstarAucP}{<0.001}
\newcommand{\FstarAucNPaired}{192}
\newcommand{\FstarAucMatched}{0.599}
\newcommand{\FstarAucMatchedBase}{0.793}
\newcommand{\FstarAucVerdict}{separated}
\newcommand{\FstarAucNS}{~\textbf{sig}}
\newcommand{\UltraAucD}{-0.1408}
\newcommand{\UltraAucDCI}{[-0.1591, -0.1224]}
\newcommand{\UltraAucP}{<0.001}
\newcommand{\UltraAucNPaired}{192}
\newcommand{\UltraAucMatched}{0.653}
\newcommand{\UltraAucMatchedBase}{0.793}
\newcommand{\UltraAucVerdict}{separated}
\newcommand{\UltraAucNS}{~\textbf{sig}}
\newcommand{\RandAucD}{-0.2933}
\newcommand{\RandAucDCI}{[-0.3135, -0.2720]}
\newcommand{\RandAucP}{<0.001}
\newcommand{\RandAucNPaired}{192}
\newcommand{\RandAucMatched}{0.500}
\newcommand{\RandAucMatchedBase}{0.793}
\newcommand{\RandAucVerdict}{separated}
\newcommand{\RandAucNS}{~\textbf{sig}}

\newcommand{\TabPr}{0.376}
\newcommand{\TabPrCI}{[0.339, 0.411]}
\newcommand{\AlgPr}{0.333}
\newcommand{\AlgPrCI}{[0.301, 0.366]}
\newcommand{\AlgLinPr}{0.359}
\newcommand{\AlgLinPrCI}{[0.324, 0.393]}
\newcommand{\PtwoRPr}{0.358}
\newcommand{\PtwoRPrCI}{[0.322, 0.394]}
\newcommand{\GeoPr}{0.276}
\newcommand{\GeoPrCI}{[0.246, 0.306]}
\newcommand{\QlutPr}{0.312}
\newcommand{\QlutPrCI}{[0.282, 0.341]}
\newcommand{\QlutSeqPr}{0.323}
\newcommand{\QlutSeqPrCI}{[0.293, 0.353]}
\newcommand{\SstarPr}{0.255}
\newcommand{\SstarPrCI}{[0.226, 0.285]}
\newcommand{\FstarPr}{0.205}
\newcommand{\FstarPrCI}{[0.177, 0.232]}
\newcommand{\UltraPr}{0.257}
\newcommand{\UltraPrCI}{[0.227, 0.287]}
\newcommand{\RandPr}{0.093}
\newcommand{\RandPrCI}{[0.084, 0.103]}
\newcommand{\TabPrD}{+0.0177}
\newcommand{\TabPrDCI}{[-0.0094, +0.0450]}
\newcommand{\TabPrP}{0.202}
\newcommand{\TabPrNPaired}{192}
\newcommand{\TabPrMatched}{0.376}
\newcommand{\TabPrMatchedBase}{0.358}
\newcommand{\TabPrVerdict}{indistinguishable}
\newcommand{\TabPrNS}{~ns}
\newcommand{\AlgPrD}{-0.0247}
\newcommand{\AlgPrDCI}{[-0.0520, +0.0028]}
\newcommand{\AlgPrP}{0.079}
\newcommand{\AlgPrNPaired}{192}
\newcommand{\AlgPrMatched}{0.333}
\newcommand{\AlgPrMatchedBase}{0.358}
\newcommand{\AlgPrVerdict}{indistinguishable}
\newcommand{\AlgPrNS}{~ns}
\newcommand{\AlgLinPrD}{+0.0014}
\newcommand{\AlgLinPrDCI}{[-0.0273, +0.0300]}
\newcommand{\AlgLinPrP}{0.926}
\newcommand{\AlgLinPrNPaired}{192}
\newcommand{\AlgLinPrMatched}{0.359}
\newcommand{\AlgLinPrMatchedBase}{0.358}
\newcommand{\AlgLinPrVerdict}{indistinguishable}
\newcommand{\AlgLinPrNS}{~ns}
\newcommand{\GeoPrD}{-0.0816}
\newcommand{\GeoPrDCI}{[-0.1087, -0.0549]}
\newcommand{\GeoPrP}{<0.001}
\newcommand{\GeoPrNPaired}{192}
\newcommand{\GeoPrMatched}{0.276}
\newcommand{\GeoPrMatchedBase}{0.358}
\newcommand{\GeoPrVerdict}{separated}
\newcommand{\GeoPrNS}{~\textbf{sig}}
\newcommand{\QlutPrD}{-0.0461}
\newcommand{\QlutPrDCI}{[-0.0740, -0.0187]}
\newcommand{\QlutPrP}{0.002}
\newcommand{\QlutPrNPaired}{192}
\newcommand{\QlutPrMatched}{0.312}
\newcommand{\QlutPrMatchedBase}{0.358}
\newcommand{\QlutPrVerdict}{separated}
\newcommand{\QlutPrNS}{~\textbf{sig}}
\newcommand{\QlutSeqPrD}{-0.0348}
\newcommand{\QlutSeqPrDCI}{[-0.0622, -0.0079]}
\newcommand{\QlutSeqPrP}{0.011}
\newcommand{\QlutSeqPrNPaired}{192}
\newcommand{\QlutSeqPrMatched}{0.323}
\newcommand{\QlutSeqPrMatchedBase}{0.358}
\newcommand{\QlutSeqPrVerdict}{separated}
\newcommand{\QlutSeqPrNS}{~\textbf{sig}}
\newcommand{\SstarPrD}{-0.1028}
\newcommand{\SstarPrDCI}{[-0.1309, -0.0742]}
\newcommand{\SstarPrP}{<0.001}
\newcommand{\SstarPrNPaired}{192}
\newcommand{\SstarPrMatched}{0.255}
\newcommand{\SstarPrMatchedBase}{0.358}
\newcommand{\SstarPrVerdict}{separated}
\newcommand{\SstarPrNS}{~\textbf{sig}}
\newcommand{\FstarPrD}{-0.1534}
\newcommand{\FstarPrDCI}{[-0.1841, -0.1221]}
\newcommand{\FstarPrP}{<0.001}
\newcommand{\FstarPrNPaired}{192}
\newcommand{\FstarPrMatched}{0.205}
\newcommand{\FstarPrMatchedBase}{0.358}
\newcommand{\FstarPrVerdict}{separated}
\newcommand{\FstarPrNS}{~\textbf{sig}}
\newcommand{\UltraPrD}{-0.1011}
\newcommand{\UltraPrDCI}{[-0.1290, -0.0738]}
\newcommand{\UltraPrP}{<0.001}
\newcommand{\UltraPrNPaired}{192}
\newcommand{\UltraPrMatched}{0.257}
\newcommand{\UltraPrMatchedBase}{0.358}
\newcommand{\UltraPrVerdict}{separated}
\newcommand{\UltraPrNS}{~\textbf{sig}}
\newcommand{\RandPrD}{-0.2647}
\newcommand{\RandPrDCI}{[-0.2980, -0.2307]}
\newcommand{\RandPrP}{<0.001}
\newcommand{\RandPrNPaired}{192}
\newcommand{\RandPrMatched}{0.093}
\newcommand{\RandPrMatchedBase}{0.358}
\newcommand{\RandPrVerdict}{separated}
\newcommand{\RandPrNS}{~\textbf{sig}}

\newcommand{\TabMcc}{0.304}
\newcommand{\TabMccCI}{[0.270, 0.338]}
\newcommand{\AlgMcc}{0.273}
\newcommand{\AlgMccCI}{[0.241, 0.304]}
\newcommand{\AlgLinMcc}{0.290}
\newcommand{\AlgLinMccCI}{[0.257, 0.321]}
\newcommand{\PtwoRMcc}{0.285}
\newcommand{\PtwoRMccCI}{[0.252, 0.316]}
\newcommand{\GeoMcc}{0.239}
\newcommand{\GeoMccCI}{[0.209, 0.268]}
\newcommand{\QlutMcc}{0.242}
\newcommand{\QlutMccCI}{[0.215, 0.270]}
\newcommand{\QlutSeqMcc}{0.245}
\newcommand{\QlutSeqMccCI}{[0.216, 0.274]}
\newcommand{\SstarMcc}{0.218}
\newcommand{\SstarMccCI}{[0.190, 0.246]}
\newcommand{\FstarMcc}{0.137}
\newcommand{\FstarMccCI}{[0.112, 0.161]}
\newcommand{\UltraMcc}{0.215}
\newcommand{\UltraMccCI}{[0.187, 0.242]}
\newcommand{\RandMcc}{n/a}
\newcommand{\RandMccCI}{[n/a, n/a]}
\newcommand{\TabMccD}{+0.0247}
\newcommand{\TabMccDCI}{[+0.0004, +0.0494]}
\newcommand{\TabMccP}{0.046}
\newcommand{\TabMccNPaired}{186}
\newcommand{\TabMccMatched}{0.309}
\newcommand{\TabMccMatchedBase}{0.285}
\newcommand{\TabMccVerdict}{separated}
\newcommand{\TabMccNS}{~\textbf{sig}}
\newcommand{\AlgMccD}{-0.0078}
\newcommand{\AlgMccDCI}{[-0.0331, +0.0171]}
\newcommand{\AlgMccP}{0.533}
\newcommand{\AlgMccNPaired}{186}
\newcommand{\AlgMccMatched}{0.277}
\newcommand{\AlgMccMatchedBase}{0.285}
\newcommand{\AlgMccVerdict}{indistinguishable}
\newcommand{\AlgMccNS}{~ns}
\newcommand{\AlgLinMccD}{+0.0071}
\newcommand{\AlgLinMccDCI}{[-0.0200, +0.0340]}
\newcommand{\AlgLinMccP}{0.603}
\newcommand{\AlgLinMccNPaired}{186}
\newcommand{\AlgLinMccMatched}{0.292}
\newcommand{\AlgLinMccMatchedBase}{0.285}
\newcommand{\AlgLinMccVerdict}{indistinguishable}
\newcommand{\AlgLinMccNS}{~ns}
\newcommand{\GeoMccD}{-0.0450}
\newcommand{\GeoMccDCI}{[-0.0695, -0.0213]}
\newcommand{\GeoMccP}{<0.001}
\newcommand{\GeoMccNPaired}{186}
\newcommand{\GeoMccMatched}{0.240}
\newcommand{\GeoMccMatchedBase}{0.285}
\newcommand{\GeoMccVerdict}{separated}
\newcommand{\GeoMccNS}{~\textbf{sig}}
\newcommand{\QlutMccD}{-0.0426}
\newcommand{\QlutMccDCI}{[-0.0662, -0.0187]}
\newcommand{\QlutMccP}{<0.001}
\newcommand{\QlutMccNPaired}{186}
\newcommand{\QlutMccMatched}{0.242}
\newcommand{\QlutMccMatchedBase}{0.285}
\newcommand{\QlutMccVerdict}{separated}
\newcommand{\QlutMccNS}{~\textbf{sig}}
\newcommand{\QlutSeqMccD}{-0.0399}
\newcommand{\QlutSeqMccDCI}{[-0.0647, -0.0147]}
\newcommand{\QlutSeqMccP}{0.003}
\newcommand{\QlutSeqMccNPaired}{186}
\newcommand{\QlutSeqMccMatched}{0.245}
\newcommand{\QlutSeqMccMatchedBase}{0.285}
\newcommand{\QlutSeqMccVerdict}{separated}
\newcommand{\QlutSeqMccNS}{~\textbf{sig}}
\newcommand{\SstarMccD}{-0.0681}
\newcommand{\SstarMccDCI}{[-0.0944, -0.0417]}
\newcommand{\SstarMccP}{<0.001}
\newcommand{\SstarMccNPaired}{186}
\newcommand{\SstarMccMatched}{0.217}
\newcommand{\SstarMccMatchedBase}{0.285}
\newcommand{\SstarMccVerdict}{separated}
\newcommand{\SstarMccNS}{~\textbf{sig}}
\newcommand{\FstarMccD}{-0.1528}
\newcommand{\FstarMccDCI}{[-0.1829, -0.1231]}
\newcommand{\FstarMccP}{<0.001}
\newcommand{\FstarMccNPaired}{186}
\newcommand{\FstarMccMatched}{0.132}
\newcommand{\FstarMccMatchedBase}{0.285}
\newcommand{\FstarMccVerdict}{separated}
\newcommand{\FstarMccNS}{~\textbf{sig}}
\newcommand{\UltraMccD}{-0.0709}
\newcommand{\UltraMccDCI}{[-0.0961, -0.0451]}
\newcommand{\UltraMccP}{<0.001}
\newcommand{\UltraMccNPaired}{186}
\newcommand{\UltraMccMatched}{0.214}
\newcommand{\UltraMccMatchedBase}{0.285}
\newcommand{\UltraMccVerdict}{separated}
\newcommand{\UltraMccNS}{~\textbf{sig}}
\newcommand{\RandMccD}{n/a}
\newcommand{\RandMccDCI}{[n/a, n/a]}
\newcommand{\RandMccP}{n/a}
\newcommand{\RandMccNPaired}{0}
\newcommand{\RandMccMatched}{n/a}
\newcommand{\RandMccMatchedBase}{n/a}
\newcommand{\RandMccVerdict}{n/a}
\newcommand{\RandMccNS}{}

\newcommand{\TabFOne}{0.333}
\newcommand{\TabFOneCI}{[0.301, 0.365]}
\newcommand{\AlgFOne}{0.307}
\newcommand{\AlgFOneCI}{[0.277, 0.337]}
\newcommand{\AlgLinFOne}{0.321}
\newcommand{\AlgLinFOneCI}{[0.290, 0.350]}
\newcommand{\PtwoRFOne}{0.302}
\newcommand{\PtwoRFOneCI}{[0.271, 0.332]}
\newcommand{\GeoFOne}{0.287}
\newcommand{\GeoFOneCI}{[0.258, 0.314]}
\newcommand{\QlutFOne}{0.272}
\newcommand{\QlutFOneCI}{[0.247, 0.298]}
\newcommand{\QlutSeqFOne}{0.275}
\newcommand{\QlutSeqFOneCI}{[0.248, 0.301]}
\newcommand{\SstarFOne}{0.267}
\newcommand{\SstarFOneCI}{[0.239, 0.293]}
\newcommand{\FstarFOne}{0.203}
\newcommand{\FstarFOneCI}{[0.179, 0.226]}
\newcommand{\UltraFOne}{0.264}
\newcommand{\UltraFOneCI}{[0.237, 0.290]}
\newcommand{\RandFOne}{0.000}
\newcommand{\RandFOneCI}{[0.000, 0.000]}
\newcommand{\TabFOneD}{+0.0315}
\newcommand{\TabFOneDCI}{[+0.0089, +0.0536]}
\newcommand{\TabFOneP}{0.006}
\newcommand{\TabFOneNPaired}{192}
\newcommand{\TabFOneMatched}{0.333}
\newcommand{\TabFOneMatchedBase}{0.302}
\newcommand{\TabFOneVerdict}{separated}
\newcommand{\TabFOneNS}{~\textbf{sig}}
\newcommand{\AlgFOneD}{+0.0048}
\newcommand{\AlgFOneDCI}{[-0.0179, +0.0269]}
\newcommand{\AlgFOneP}{0.692}
\newcommand{\AlgFOneNPaired}{192}
\newcommand{\AlgFOneMatched}{0.307}
\newcommand{\AlgFOneMatchedBase}{0.302}
\newcommand{\AlgFOneVerdict}{indistinguishable}
\newcommand{\AlgFOneNS}{~ns}
\newcommand{\AlgLinFOneD}{+0.0186}
\newcommand{\AlgLinFOneDCI}{[-0.0066, +0.0430]}
\newcommand{\AlgLinFOneP}{0.149}
\newcommand{\AlgLinFOneNPaired}{192}
\newcommand{\AlgLinFOneMatched}{0.321}
\newcommand{\AlgLinFOneMatchedBase}{0.302}
\newcommand{\AlgLinFOneVerdict}{indistinguishable}
\newcommand{\AlgLinFOneNS}{~ns}
\newcommand{\GeoFOneD}{-0.0154}
\newcommand{\GeoFOneDCI}{[-0.0381, +0.0073]}
\newcommand{\GeoFOneP}{0.183}
\newcommand{\GeoFOneNPaired}{192}
\newcommand{\GeoFOneMatched}{0.287}
\newcommand{\GeoFOneMatchedBase}{0.302}
\newcommand{\GeoFOneVerdict}{indistinguishable}
\newcommand{\GeoFOneNS}{~ns}
\newcommand{\QlutFOneD}{-0.0295}
\newcommand{\QlutFOneDCI}{[-0.0521, -0.0073]}
\newcommand{\QlutFOneP}{0.011}
\newcommand{\QlutFOneNPaired}{192}
\newcommand{\QlutFOneMatched}{0.272}
\newcommand{\QlutFOneMatchedBase}{0.302}
\newcommand{\QlutFOneVerdict}{separated}
\newcommand{\QlutFOneNS}{~\textbf{sig}}
\newcommand{\QlutSeqFOneD}{-0.0270}
\newcommand{\QlutSeqFOneDCI}{[-0.0507, -0.0033]}
\newcommand{\QlutSeqFOneP}{0.023}
\newcommand{\QlutSeqFOneNPaired}{192}
\newcommand{\QlutSeqFOneMatched}{0.275}
\newcommand{\QlutSeqFOneMatchedBase}{0.302}
\newcommand{\QlutSeqFOneVerdict}{separated}
\newcommand{\QlutSeqFOneNS}{~\textbf{sig}}
\newcommand{\SstarFOneD}{-0.0353}
\newcommand{\SstarFOneDCI}{[-0.0597, -0.0100]}
\newcommand{\SstarFOneP}{0.006}
\newcommand{\SstarFOneNPaired}{192}
\newcommand{\SstarFOneMatched}{0.267}
\newcommand{\SstarFOneMatchedBase}{0.302}
\newcommand{\SstarFOneVerdict}{separated}
\newcommand{\SstarFOneNS}{~\textbf{sig}}
\newcommand{\FstarFOneD}{-0.0993}
\newcommand{\FstarFOneDCI}{[-0.1271, -0.0714]}
\newcommand{\FstarFOneP}{<0.001}
\newcommand{\FstarFOneNPaired}{192}
\newcommand{\FstarFOneMatched}{0.203}
\newcommand{\FstarFOneMatchedBase}{0.302}
\newcommand{\FstarFOneVerdict}{separated}
\newcommand{\FstarFOneNS}{~\textbf{sig}}
\newcommand{\UltraFOneD}{-0.0378}
\newcommand{\UltraFOneDCI}{[-0.0622, -0.0128]}
\newcommand{\UltraFOneP}{0.002}
\newcommand{\UltraFOneNPaired}{192}
\newcommand{\UltraFOneMatched}{0.264}
\newcommand{\UltraFOneMatchedBase}{0.302}
\newcommand{\UltraFOneVerdict}{separated}
\newcommand{\UltraFOneNS}{~\textbf{sig}}
\newcommand{\RandFOneD}{-0.3019}
\newcommand{\RandFOneDCI}{[-0.3318, -0.2713]}
\newcommand{\RandFOneP}{<0.001}
\newcommand{\RandFOneNPaired}{192}
\newcommand{\RandFOneMatched}{0.000}
\newcommand{\RandFOneMatchedBase}{0.302}
\newcommand{\RandFOneVerdict}{separated}
\newcommand{\RandFOneNS}{~\textbf{sig}}



\title{Algebraic and topological invariant fields for cryptic-pocket detection\\
on the official CryptoBench receptor-only benchmark}
\author{Shane \and Xuening Wu}
\date{July 2026}

\begin{document}
\maketitle

\begin{abstract}
We evaluate deterministic, training-free geometric detectors on the task of
recovering cryptic binding-site residues from apo receptors, using the official
CryptoBench test fold (MMseqs2 clustering at 10\% sequence identity,
cluster-disjoint). The detectors have no trained weights and no fitted
hyperparameters, and were never exposed to any CryptoBench partition; the input
contract is receptor coordinates only, enforced at runtime by a ligand-leak
guard that rejects any non-solvent heteroatom. On \NUnits{} single-chain
evaluation units a rigid geometric detector with no training of any kind attains
a per-residue ROC-AUC of \GeoAuc{} (95\% CI \GeoAucCI{}), while a locally
executed P2Rank 2.5.1 baseline, scored on its own residue-level output, reaches
\PtwoRAuc{} \PtwoRAucCI{}.

Between the two we place detectors of a third kind, built entirely from exact
algebraic and topological invariants of the receptor and compiled on the
cluster-disjoint CryptoBench training fold. Each invariant is banded into a
quaternary digit against its own chain's order statistics; the detectors differ
only in how those digits become a score. Reading them through dense integer
counter tables fused by integer-multiplicity addition gives a field whose entire
resolution path is counting: it attains ROC-AUC \AlgAuc{} \AlgAucCI{} and is
separated from P2Rank on ROC-AUC (\AlgAucD{} \AlgAucDCI{}). Reading the same
digits through a single regularised closed-form linear solve --- one symmetric
system, no gradient, no iteration, no auto-differentiation --- gives
ROC-AUC \AlgLinAuc{} \AlgLinAucCI{}, PR-AUC \AlgLinPr{} \AlgLinPrCI{}, MCC
\AlgLinMcc{} \AlgLinMccCI{} and F1 \AlgLinFOne{} \AlgLinFOneCI{}. Against P2Rank
on the same resample indices its paired differences are \AlgLinAucD{}
\AlgLinAucDCI{} on ROC-AUC, \AlgLinPrD{} \AlgLinPrDCI{} on PR-AUC, \AlgLinMccD{}
\AlgLinMccDCI{} on MCC and \AlgLinFOneD{} \AlgLinFOneDCI{} on F1: all four
intervals contain zero. We claim parity with the trained baseline on all four
metrics, and nothing stronger --- the point estimates favour the field on all
four, but not by margins this fold can resolve.

The gap between the two readouts is a sample-size effect and we quantify it
rather than assert it. A width-six quaternary table has \(4^{6}\) cells, so a
cell holds about fifty training residues and three positives and its fraction
carries a relative standard error near 60\%; a first-order functional estimates
each coefficient from the whole fold. Varying the table width shows the
trade-off directly: width-two tables, whose cells are estimated to 3\%, recover
second-order structure but score below the first-order functional, so on this
fold the interactions available to a quaternary code do not pay for their
variance.

Both the readout family and the fusion architecture were selected on a
cluster-disjoint split of the training fold alone, before the test fold was
read, so no reported figure is the maximum of several test-fold estimates. Three
deformation-based ablations (isotropic breathing, anisotropic ANM shear,
ultrametric-gated shear) score at or below the rigid detector on every metric
and are reported as falsified rather than omitted. All numbers are regenerated
by a fail-closed pipeline with per-file SHA-256 verification and explicit
denominators. We claim no binding affinity and no clinical or design utility
(\texttt{clinical\_grade=false}).
\end{abstract}

\section{Scope}
This paper makes two claims about the official CryptoBench receptor-only test
fold. First, a detector built only from exact algebraic and topological
invariants of the apo receptor, read out by a single closed-form linear solve,
is not separable from a locally executed P2Rank on any of ROC-AUC, PR-AUC, MCC
or F1. Second, restricting the same invariants to a resolution path of pure
integer counting costs \(\AlgLinAuc{}-\AlgAuc{}\) ROC-AUC on this fold, and we
attribute that cost to the variance of a cell fraction rather than to any
deficiency of the invariants.

Explicitly not claimed: superiority over trained baselines, parity on any other
benchmark, binding affinity, candidate efficacy, drug design utility, clinical
readiness, or global novelty. Both fields are compiled on the CryptoBench
training fold and are therefore \emph{not} training-free, and the linear field
additionally fits real-valued coefficients; the training-free detectors are
reported separately and are weaker. Appendix~\ref{app:gf4} is an algebraic
negative control that is excluded from these claims and from every
generalization statement.

\section{Data}
The benchmark is the official CryptoBench test fold (Skrhak et al.,
\emph{Bioinformatics} 41(1):btae745, 2025; dataset \url{https://osf.io/pz4a9/}),
whose splits are cluster-disjoint under MMseqs2 at 10\% sequence identity.
Receptors are materialized from RCSB and every receptor and label file is pinned
by SHA-256 in a manifest; the loader refuses to run if the manifest is absent,
if the declared clustering method or identity threshold differs, if any hash
mismatches, or if a cluster identifier appears in more than one split. There is
no fallback path: a smaller pilot subset cannot be substituted for this fold.

The evaluation unit is a \((\text{pdb},\text{chain})\) pair, not a PDB entry:
192 units span 190 distinct entries because two entries contribute two chains
each. A further 38 fold units are multi-chain assemblies and are excluded before
evaluation, enumerated in the manifest, because the residue indexing used here is
chain-agnostic and cannot address such assemblies unambiguously. The exclusion is
recorded rather than silent.

\section{Methods}
\paragraph{Input contract.} Every detector receives receptor coordinates only.
A decorator asserts, at call time, that the file carries no non-solvent HETATM
record and no non-polymer ATOM record; a violation is recorded as a crash, not
skipped. An import-graph test additionally forbids any detector module from
importing the scorer, so a label cannot reach a predictor through code.

\paragraph{Rigid geometric detector.} Free grid points outside the van der Waals
surface are scored by a buriedness functional, the fraction of probe directions
blocked by an atom within a ray cylinder. Candidate centres are the most buried
points, deduplicated by a minimum separation. There are no trained weights and no
fitted hyperparameters; the geometric constants are fixed a priori and are not
tuned against any CryptoBench partition.

\paragraph{Deformation ablations.} Two ablations test whether admitting
deformation improves recovery of pockets that are closed in the apo frame. The
isotropic ablation applies a fixed set of radial dilations. The anisotropic
ablation replaces isotropic dilation with the receptor's own low-frequency,
non-rigid modes from an anisotropic network model: a candidate is promoted when a
small shear opens a large still-enclosed cavity at that site. Shear amplitudes
are fixed by the spectrum as \(c_k=\Delta x\sqrt{\lambda_1/\lambda_k}\), so no
scalar is tuned. The eigen-decomposition is pinned for determinism by a fixed
Lanczos start vector, zero convergence tolerance, and an eigenvector phase
convention; repeated runs are byte-identical.

\paragraph{Three classes of detector, and what ``training-free'' means here.}
The detectors above see no CryptoBench partition of any kind: their constants are
fixed a priori and the apo receptor is their only input. The field detectors
below are not of that class and we do not describe them as training-free. They
are compiled on the official CryptoBench training fold, which is cluster-disjoint
from the test fold at 10\% sequence identity under the same MMseqs2 partition.
Among them we draw one further distinction, and it is mechanical rather than
rhetorical. The counting fields store only pairs of integers --- how many
residues addressed a cell, and how many of them were cryptic --- and every
quantity they read at inference time is a ratio of two such integers; no gradient
step is taken, no real-valued weight is solved for, and no continuous
discriminant is fitted. The linear field relaxes exactly that last property: it
solves one regularised symmetric system for a real coefficient vector. It still
takes no gradient step, runs no iteration, differentiates nothing, and reads no
test residue, but it is a fitted model and every table in this paper labels it as
one. We separate the three classes rather than averaging over them, because
averaging over what a method was allowed to see is what makes such tables
misleading.

\paragraph{Quaternary resolution field.} Six geometric invariants per residue are
banded into quaternary digits against the residue's own chain quartiles, which
is a comparator network over the chain and carries no constant between chains.
The digit word addresses a dense table of \(4^{6}\) cells holding the two
integer counters. The addressable width is not free: with \(N\) training residues
at cryptic base rate \(r\), a cell holds \(rN/4^{d}\) expected positives, and
requiring at least one gives \(d \le \log_4(rN)\). On this fold
\(rN = 13\,524\) and the bound is \(d \le 6.87\), so a flat word may not exceed
seven digits. An earlier ten-digit variant violated it: \(4^{10}\) cells against
\(13\,524\) positives left 98.9\% of the address space undrivable.

\paragraph{Algebraic counting field.} To reach the descriptive richness of a
feature-based predictor without leaving exact quantities, the word is widened to
35 invariants, each a closed-form algebraic or integer topological function of
the receptor: solid-angle occlusion and protrusion; the Laplacian spectrum of the
induced neighbourhood graph together with the Betti numbers of the local pocket
lining; the closed-form least-squares gradient and discrete Laplacian of the
packing-density field; non-Archimedean ball geometry and the residue valuation;
inertia-tensor anisotropy, planarity and sphericity; and depth, coordination and
contact order. Because the capacity bound forbids addressing 35 digits at once,
the invariants are partitioned into six thematic blocks of at most six digits,
each block compiled into its own dense table, and the blocks fused by addition of
their cell fractions. Fusing rather than cascading is deliberate: a block table
estimates its fraction from about 57 residues, whose standard error is 53\% of
the base rate, and banding such an estimate into two bits is an irreversible
decision taken on a noise-dominated statistic. Measured on this fold, every
cascaded topology we compiled scored below the flat control, while parallel
fusion scored above it. Block multiplicities in the sum are small integers given
by the rank of each block's own compiled Gini index, i.e.\ fan-out replication of
a table, not a real coefficient. A final spatial gate adds, at five radii, the
mean of the fused score over the closed geometric neighbourhood, which states
combinationally the fact that a cryptic site is a contiguous patch of residues
rather than an isolated one.

\paragraph{Wire expansion.} The linear field widens the word from the 35
invariants to \NLinWires{} wires without introducing any learned representation.
Seven published per-residue constants are appended (Kyte--Doolittle hydropathy,
side-chain volume, aromatic ring count, formal charge at pH 7, hydrogen-bond
donor and acceptor counts, and rotatable side-chain bonds), together with one
counted propensity \(P(\text{cryptic}\mid\text{residue type})\) compiled on the
training fold alone and carried into the artifact, so no test residue ever
contributes a count to it. Each of the resulting 43 quantities is then taken both
at the residue and as its mean over the 6, 14 and 20\,\AA{} neighbourhoods, which
states at three scales the fact that a cryptic site is a property of a
neighbourhood rather than of one residue. Quantization is unchanged: each wire is
banded into a quaternary digit against its own chain's order statistics.

\paragraph{Linear readout, and why table width is the free variable.} The readout
is one regularised solve, \(w=(A^{\top}A+\lambda\,\bar{\sigma}I)^{-1}A^{\top}t\),
over the standardised digits, with \(\lambda=\LinRidge{}\) and \(t\) the centred
label. The score is then \(s+\LinGateWeight{}\,\bar{s}\), where \(\bar{s}\) is
the mean of \(s\) over the closed \LinGateRadius{}\,\AA{} neighbourhood rescaled
to the spread of \(s\) --- the same patch statement the counting field makes,
applied to a continuous field.

Which readout to use is not a matter of taste here, and the quantity that decides
it is the width of the table. A width-\(d\) quaternary table has \(4^{d}\) cells,
so on \(N=\NTrainResidues{}\) training residues at base rate \TrainRate{} a cell
holds \(N/4^{d}\) residues and its fraction has relative standard error
\(\sqrt{(1-r)/(rN/4^{d})}\). At \(d=6\) that is about 60\%: the bank models the
joint distribution inside each window exactly, but each exact statement is about
a noisy quantity. At \(d=2\) it is about 3\%, and there are only \(\binom{K}{2}\)
such tables over the top \(K\) wires, few enough to combine by the same single
solve. A first-order functional sits at the other end, representing no
interaction at all but estimating each coefficient from the whole fold. We
therefore compiled all three and let the training fold choose, which is reported
in Section~\ref{sec:readout} rather than assumed here.

\paragraph{Metrics.} All four metrics are computed over the same residue universe
(every residue of the apo chain) at the same operating point (residue belongs to
any predicted pocket): ROC-AUC, average precision (PR-AUC), MCC and F1.
Uncertainty is a paired bootstrap over evaluation units, 10{,}000 resamples, 95\%
percentile CI, with the same resample indices applied to every method so that
between-method correlation is preserved. A metric that is undefined for a
structure is null and is dropped from that metric's denominator; it is never
imputed as zero. Per-unit values are keyed on \((\text{pdb},\text{chain})\), and a
colliding key raises rather than overwriting a row.

\paragraph{Baselines.} P2Rank 2.5.1 is executed locally with a pinned version and
is a scored baseline. A uniform bounding-box sampler provides a chance floor.
PocketMiner is reported as data-unavailable: the published artifact is a trained
model binary with no per-residue predictions, so no paired interval against it is
computed and no value is imputed. Tool unavailability is never counted as a miss.

\section{Results}
The official evaluation emits \NRows{} telemetry rows, \NUnits{} evaluation
units times \NMethods{} methods, with per-residue metrics available on all of
them and no silent drops. Every method returns a prediction on
\NUnits{}/\NUnits{} units.

P2Rank is scored on the residue-level table it natively emits
(\texttt{*\_residues.csv}): the calibrated \texttt{probability} column as the
continuous per-residue score, and its own \texttt{pocket > 0} assignment as the
binary operating point. An earlier revision of this harness instead reconstructed
a residue signal from pocket centres, keeping the top 5 pockets and weighting
every residue within 6\,\AA{} of a centre by \(1/\text{rank}\). That is not
P2Rank's prediction and it is lossy: \(1/\text{rank}\) is a five-valued staircase
that cannot resolve the residue ordering P2Rank emits, and truncation ties every
remaining residue at zero. Under that harness P2Rank measured ROC-AUC 0.621;
scored on its own output it reaches \PtwoRAuc{}, consistent with the published
CryptoBench baseline. We report the corrected figure.

\begin{table}[h]
\centering
\small
\setlength{\tabcolsep}{2pt}
\begin{tabular}{@{}lllll@{}}
\toprule
Method & ROC-AUC & PR-AUC & MCC & F1\\
\midrule
\multicolumn{5}{@{}l}{\emph{trained on CryptoBench}}\\
P2Rank 2.5.1          & \PtwoRAuc{} \PtwoRAucCI{} & \PtwoRPr{} \PtwoRPrCI{}
                      & \PtwoRMcc{} \PtwoRMccCI{} & \PtwoRFOne{} \PtwoRFOneCI{}\\
\addlinespace
\multicolumn{5}{@{}l}{\emph{fitted on the training fold (one closed-form solve)}}\\
Algebraic field (linear) & \AlgLinAuc{} \AlgLinAucCI{} & \AlgLinPr{} \AlgLinPrCI{}
                      & \AlgLinMcc{} \AlgLinMccCI{} & \AlgLinFOne{} \AlgLinFOneCI{}\\
\addlinespace
\multicolumn{5}{@{}l}{\emph{compiled by counting on the training fold}}\\
Algebraic field       & \AlgAuc{} \AlgAucCI{} & \AlgPr{} \AlgPrCI{}
                      & \AlgMcc{} \AlgMccCI{} & \AlgFOne{} \AlgFOneCI{}\\
Quaternary field (+seq) & \QlutSeqAuc{} \QlutSeqAucCI{} & \QlutSeqPr{} \QlutSeqPrCI{}
                      & \QlutSeqMcc{} \QlutSeqMccCI{} & \QlutSeqFOne{} \QlutSeqFOneCI{}\\
Quaternary field      & \QlutAuc{} \QlutAucCI{} & \QlutPr{} \QlutPrCI{}
                      & \QlutMcc{} \QlutMccCI{} & \QlutFOne{} \QlutFOneCI{}\\
\addlinespace
\multicolumn{5}{@{}l}{\emph{training-free}}\\
Rigid geometric       & \GeoAuc{} \GeoAucCI{} & \GeoPr{} \GeoPrCI{}
                      & \GeoMcc{} \GeoMccCI{} & \GeoFOne{} \GeoFOneCI{}\\
Anisotropic shear     & \SstarAuc{} \SstarAucCI{} & \SstarPr{} \SstarPrCI{}
                      & \SstarMcc{} \SstarMccCI{} & \SstarFOne{} \SstarFOneCI{}\\
Ultrametric shear     & \UltraAuc{} \UltraAucCI{} & \UltraPr{} \UltraPrCI{}
                      & \UltraMcc{} \UltraMccCI{} & \UltraFOne{} \UltraFOneCI{}\\
Isotropic ablation    & \FstarAuc{} \FstarAucCI{} & \FstarPr{} \FstarPrCI{}
                      & \FstarMcc{} \FstarMccCI{} & \FstarFOne{} \FstarFOneCI{}\\
Chance (bbox)         & \RandAuc{} \RandAucCI{} & \RandPr{} \RandPrCI{}
                      & \RandMcc{} & \RandFOne{} \RandFOneCI{}\\
\bottomrule
\end{tabular}
\caption{Per-residue metrics on the official CryptoBench test fold, point
estimate with 95\% paired-bootstrap CI over \(n=\NUnits{}\) evaluation units,
\NBoot{} resamples, seed \BootSeed{}. Methods are grouped by what they were
allowed to see, because averaging over that distinction is what makes such
tables misleading. A metric that is undefined on a unit is null and leaves that
unit out of its own denominator; the chance baseline's MCC is undefined on every
unit because no residue crosses the operating point, and is reported as such
rather than imputed as zero.}
\end{table}

\begin{table}[h]
\centering
\small
\setlength{\tabcolsep}{3pt}
\begin{tabular}{@{}lll@{}}
\toprule
Method \(-\) P2Rank & ROC-AUC & PR-AUC\\
\midrule
Algebraic field (linear) & \AlgLinAucD{} \AlgLinAucDCI{}\AlgLinAucNS{}
                         & \AlgLinPrD{} \AlgLinPrDCI{}\AlgLinPrNS{}\\
Algebraic field          & \AlgAucD{} \AlgAucDCI{}\AlgAucNS{}
                         & \AlgPrD{} \AlgPrDCI{}\AlgPrNS{}\\
Quaternary field         & \QlutAucD{} \QlutAucDCI{}\QlutAucNS{}
                         & \QlutPrD{} \QlutPrDCI{}\QlutPrNS{}\\
Rigid geometric          & \GeoAucD{} \GeoAucDCI{}\GeoAucNS{}
                         & \GeoPrD{} \GeoPrDCI{}\GeoPrNS{}\\
\midrule
Method \(-\) P2Rank & MCC & F1\\
\midrule
Algebraic field (linear) & \AlgLinMccD{} \AlgLinMccDCI{}\AlgLinMccNS{}
                         & \AlgLinFOneD{} \AlgLinFOneDCI{}\AlgLinFOneNS{}\\
Algebraic field          & \AlgMccD{} \AlgMccDCI{}\AlgMccNS{}
                         & \AlgFOneD{} \AlgFOneDCI{}\AlgFOneNS{}\\
Quaternary field         & \QlutMccD{} \QlutMccDCI{}\QlutMccNS{}
                         & \QlutFOneD{} \QlutFOneDCI{}\QlutFOneNS{}\\
Rigid geometric          & \GeoMccD{} \GeoMccDCI{}\GeoMccNS{}
                         & \GeoFOneD{} \GeoFOneDCI{}\GeoFOneNS{}\\
\bottomrule
\end{tabular}
\caption{Paired differences against the locally executed P2Rank baseline on the
same resample indices, so the within-unit correlation between two detectors
looking at the same pocket is removed from the difference rather than inflating
its variance. ``Indistinguishable'' means the 95\% interval of the paired
difference contains zero. With the linear readout all four intervals contain
zero and all four point estimates are positive; with the counting readout the
fold separates the field from P2Rank on ROC-AUC (\AlgAucD{},
\(p=\)~\AlgAucP{}). A positive point estimate inside an interval that contains
zero is not evidence of superiority, and we do not read it as such.}
\end{table}

\paragraph{The architecture was not chosen on the test fold.} \NCandidates{}
fusion architectures were compared during development. Selecting among them by
test-fold score and then reporting the winner's test score would make the
reported figure the maximum of \NCandidates{} noisy estimates. The choice was
therefore made on the training fold alone, split into two halves with disjoint
MMseqs2 clusters (\NSelectFitUnits{} units to compile on, \NSelectPickUnits{} to
rank on), and the official test fold is read once, for the single architecture
that procedure names. It names ``\SelectedArch{}'' at ROC-AUC
\SelectedArchAuc{} on the held-out training half, which is the architecture
reported here; the training-only ranking of all \NCandidates{} candidates agrees
with the test-fold ranking, so the test figure can be read as an out-of-sample
estimate rather than a selected maximum.

\paragraph{What limits the counting field.} Holding the estimator fixed and
widening the invariant word from six to \NAlgFeatures{} moves a closed-form
linear functional of the same invariants from ROC-AUC \FisherSix{} to \FisherAll{} on this
fold, so the descriptors carry the signal. Pushed through counter tables the
same invariants reach \AlgAuc{}. The gap is a capacity limit rather than a
descriptive one: a table over \(d\) quaternary digits has \(4^{d}\) free cells
and is bounded by \(d \le \log_4(rN) = 6.87\) here, whereas a linear form over
\(M\) rank features is determined by \(\Theta(M^{2})\) quantities, which for
\(M=\NAlgFeatures{}\) is \(1\,225\) against \(13\,524\) available positives. The
linear form can therefore address all \NAlgFeatures{} dimensions at once while no
admissible table can address more than seven, and unweighted or
integer-weighted partitions of the invariants do not close that gap. We report
this as the measured ceiling of the counting construction, not as a conjecture.

\subsection{Choosing the readout on the training fold}
\label{sec:readout}
The readout was chosen by the same discipline as the fusion architecture and on
the same cluster-disjoint halves of the training fold: compile on
\NSelectFitUnits{} units, rank on \NSelectPickUnits{} units whose clusters the
first half never saw, and read the test fold once for the winner. On that
held-out training half the integer-multiplicity table bank reaches \PickCount{},
a linear functional of the \NLinWires{} first-order digits reaches
\PickLinDigits{}, and the same functional over the digit \emph{indicators} ---
which can represent a non-monotone response to a wire, at four times the
parameter count --- reaches \PickLinOnehot{}, slightly worse. A linear
combination of the bank's own cell fractions reaches only \PickLinTables{}: the
tables lose more to cell variance than they gain from modelling the joint
distribution, and no way of weighting them recovers it.

That diagnosis predicts where an interaction model should live, so we tested it
directly by varying the table width. Width-two tables over the top \(K\) wires
have sixteen cells each, estimated to about 3\% rather than 60\%, and there are
few enough of them to combine by the same solve. Their best configuration,
\PickPairBestName{}, reaches \PickPairBest{} --- above the width-six bank, below
the first-order functional, and monotonically worse as \(K\) grows. The
conclusion this fold supports is narrow and worth stating precisely: pairwise
interactions among quaternary-coded algebraic invariants are estimable here, and
they carry no signal beyond first order that survives their own variance.

One methodological note belongs in the record. A first pass at this sweep
accumulated the \(1770\times1770\) Gram matrix in single precision and reported
that adding tables destroys the fit, collapsing from 0.74 to 0.60. Pair
fractions are strongly collinear --- two tables sharing a wire agree on a quarter
of their address --- so the Gram matrix is near-singular and single-precision
accumulation loses exactly the small eigenvalues that carry the interaction
signal. The collapse was round-off, not a property of the model class. We report
it because the failure mode is invisible in the output: the numbers are
plausible, ordered, and wrong.

\paragraph{Negative result: deformation does not help.} Both deformation
ablations score strictly below the rigid detector on all four metrics. The
anisotropic modes used here are global and low-frequency; they act as
macroscopic hinges that separate surface regions and create decoy volume, which
drags the predicted centre away from the true site. The correct next step
indicated by this failure is spatially gated, localized shear with rigid-hit
anchoring. We report this as a falsified hypothesis rather than removing the
ablations from the paper.

\section{Limitations}
The comparison covers one trained baseline, executed locally; PocketMiner could
not be included because per-residue predictions are not published. The 38
excluded multi-chain units mean the reported fold is a strict subset of the
official test fold, and the result does not transfer to multi-chain assemblies
without a chain-aware residue index.

All four paired intervals between the linear field and P2Rank contain zero,
which licenses ``we cannot separate these on this fold'' and not ``these are
equal''. All four point estimates favour the field, and that is not evidence of
superiority: with \NUnits{} units the intervals are roughly \(\pm0.02\) wide on
ROC-AUC, so a margin of \AlgLinAucD{} is exactly the size of difference this fold
cannot resolve. A reader who wants an ordering between these two methods will
not get one here, and should not infer one from the sign. Nor does parity on this
benchmark imply parity in general: one trained baseline, executed locally, on one
receptor-only task.

The absolute values remain modest --- ROC-AUC \AlgLinAuc{} and F1 \AlgLinFOne{}
for the strongest field here --- so the contribution is the construction rather
than the accuracy. What the construction buys is auditability: every input is a
closed-form algebraic or integer topological function of the receptor, the entire
fitted state is \NLinWires{} coefficients from one symmetric solve, and the
counting variant reaches \AlgAuc{} while storing nothing but integer pairs.
Whether that is worth anything depends on the application and not on this
measurement. No affinity, design or clinical claim follows from any number in
this paper.

\section{Reproducibility}
Every number above is regenerated from the repository by two commands: the
official-fold runner (fail-closed loader, per-file SHA-256 verification) and the
paired-bootstrap module (10{,}000 resamples, fixed seed). Parallelism across
structures is a scheduling change only and is order-preserving. Optional native
kernels are an operation-for-operation port of the reference implementation and
are covered by a bit-identity test, so enabling them cannot move a reported
value. Frozen artifacts and the telemetry schema are committed alongside the
code.

% Appendix B — included by MAIN_CRYPTOBENCH_GEOAUDIT.tex via \input.
% Not a standalone article: no preamble, no \documentclass, no abstract.
%
% SCOPE BOUNDARY (machine-enforced by contracts/GEOAUDIT_PAPER_SCOPE.json and
% tools/verify_claims.py): everything below is an algebraic negative control.
% It is excluded from the primary claim and from every generalization
% statement in the main text. It contains no chemical, binding, docking,
% affinity, efficacy or clinical content, and it contributes nothing to the
% CryptoBench result reported in the main text.

\appendix
\section{Appendix B: Finite-field allele-conditioning ablation}
\label{app:gf4}

This appendix records an algebraic ablation over \(GF(4)\): nucleotide 48-mers
are mapped to finite-field tensors, an allele induces a sparse residual
\(\delta\), and a tensor \(x\) is constrained by the syndrome equation
\(Hx=B\delta\) with operators \((H,B)\) synthesized from local sequence
background. It is retained as a negative control and as future work. It is
\emph{not} evidence for the main text's benchmark result, and it asserts
nothing about affinity, steric clearance, synthesis, or clinical performance
(\texttt{clinical\_grade=false}).

\subsection{Construction}
To establish an \(O(1)\) verification environment, we deprecate the
tautological exact-form null identity
\(T(s)=s+\alpha\). Instead, we implement a non-trivial, falsifiable constrained
syndrome equation
\begin{equation}
  Hx = B\delta
\end{equation}
to isolate spatial geometric variance derived strictly from mutant alleles,
ensuring deterministic topological constraints without biological iteration.

\subsection{GF(4) Isomorphism}
We establish a strict isomorphism mapping nucleotide bases to the finite field
\(GF(4)\cong\mathbb{F}_2[x]/(x^2+x+1)\) with characteristic~2
(where \(y\oplus y=0\) and \(\alpha^2+\alpha+1=0\)). The physical base levels
are mapped as follows:
\begin{equation}
\phi(A)=0,\quad \phi(C)=1,\quad \phi(G)=\alpha,\quad \phi(U)=\alpha^2.
\end{equation}

\subsection{Sparse Residual \(\delta\) Extraction}
For the 48-mer tensor \(s\in GF(4)^{48}\) centered at codon \(k\)
(e.g., KRAS G12), we isolate the exact spatial disparity.

First, we define the local reference tensor for the wild-type (GGU):
\begin{equation}
(s_{\mathrm{ref}}^{(k)},\, s_{\mathrm{ref}}^{(k+1)},\, s_{\mathrm{ref}}^{(k+2)})
= (\alpha,\,\alpha,\,\alpha^2).
\end{equation}

Next, we define the local mutant tensors for G12C (UGU), G12D (GAU), and
G12V (GUU) in \(O(1)\) time:
\begin{align}
\text{G12C (UGU)}
&\implies
(s_{G12C}^{(k)},\, s_{G12C}^{(k+1)},\, s_{G12C}^{(k+2)})
= (\alpha^2,\,\alpha,\,\alpha^2), \\
\text{G12D (GAU)}
&\implies
(s_{G12D}^{(k)},\, s_{G12D}^{(k+1)},\, s_{G12D}^{(k+2)})
= (\alpha,\,0,\,\alpha^2), \\
\text{G12V (GUU)}
&\implies
(s_{G12V}^{(k)},\, s_{G12V}^{(k+1)},\, s_{G12V}^{(k+2)})
= (\alpha,\,\alpha^2,\,\alpha^2).
\end{align}

We compute the pure spatial deviation tensor \(\delta\in GF(4)^{48}\) via
field addition (XOR on the bit lift):
\begin{equation}
\delta = s_{\mathrm{mut}} \oplus s_{\mathrm{ref}}.
\end{equation}

By construction, all unmutated background loci annihilate instantaneously
(\(s_{\mathrm{mut}}^{(i)}\oplus s_{\mathrm{ref}}^{(i)}=0\)). The sparse non-zero
topological residuals strictly resolve to absolute field elements:
\begin{align}
\delta_{G12C}^{(k)} &= \alpha^2 \oplus \alpha = 1, \\
\delta_{G12D}^{(k+1)} &= 0 \oplus \alpha = \alpha, \\
\delta_{G12V}^{(k+1)} &= \alpha^2 \oplus \alpha = 1.
\end{align}

\subsection{Non-Trivial Inverse Solver}
The sparse residuals \(\delta_{G12C}\), \(\delta_{G12D}\), and \(\delta_{G12V}\)
represent strictly orthogonal input excitation vectors. We substitute the
exact-form logic gate with the rigorous syndrome solver:
\begin{equation}
Hx = B\delta.
\end{equation}
Where \(H\) represents the topological constraint matrix, \(B\) is the boundary
mapping operator, and \(x\) is the target solution tensor representing a
candidate molecular entity under algebraic conditioning. A geometry is treated
as algebraically admissible if and only if its tensor \(x\) deterministically
absorbs the topological tension \(\delta\) generated by the allele variance.
The solution \(x\) is conditioning evidence, not a unique chemical inverse.

\subsection{Deprecation of Tautological Complements}
The formulation \(\mathbf{S}_{\mathrm{path}}\oplus\mathbf{S}_{\mathrm{thera}}=\mathbf{1}\)
(and its associated exact-form null \(T(s)=s+\alpha\)) is mathematically
trivial: it functions as a fixed topological schedule rather than a
falsifiable chemical reverse-solution. The mandated transition is therefore to
the non-trivial syndrome solver \(Hx=B\delta\), ensuring the solution is
physically constrained by allele variance, not axiomatically forced.

\subsection{Background-Conditioned Dynamic Operators}
To remediate the pathology in which identical topological tensors \(\delta\)
yielded indiscriminate solution vectors \(x\), we abandon static operator
placeholders in favor of background-conditioned dynamic operators.

\subsubsection{Resolution of \(\delta\)-Collision via Toeplitz Modulation}
The solution space \(x\) for \(Hx=B\delta\) is not solely a function of the
mutation tensor \(\delta\), but is fundamentally constrained by the local
genomic context \(\mathbf{S}_{\mathrm{bg}}\in GF(4)^{47}\). Operator synthesis
is defined by
\begin{equation}
    H_{\mathrm{target}} = \mathcal{F}(\mathbf{S}_{\mathrm{bg}}),\qquad
    B_{\mathrm{target}} = \mathcal{G}(\mathbf{S}_{\mathrm{bg}}),
\end{equation}
where \(\mathcal{F}\) and \(\mathcal{G}\) act as Toeplitz-based spatial filters
mapping the 47-mer background residues into the operator matrices. Thus even
when two alleles (e.g., KRAS G12C and FLT3 D835Y) generate identical sparse
residuals \(\delta[23]=1\), the induced operator \(H_{\mathrm{target}}\)
diverts the solution into non-colliding null spaces.

\subsubsection{Precision Probe Pipeline}
Every probe (48-mer) is anchored to its provenance ledger via SHA-256 hashes
for reproducibility of the input tensors:
\begin{itemize}
    \item \textbf{Index-23 Anchor:} the mutation site is strictly pinned at the
    median index of the 48-mer window, ensuring algebraic alignment.
    \item \textbf{Symbolic XOR Resolution:} the spatial disparity \(\delta\) is
    computed via field addition \(\delta=s_{\mathrm{mut}}\oplus s_{\mathrm{ref}}\)
    in \(GF(4)\), isolating the non-zero topological residual to the mutation
    site.
\end{itemize}

\subsection{Truth Boundaries}
\begin{center}
\begin{tabular}{@{}lp{0.62\textwidth}@{}}
\toprule
Claim in this ledger & Status\\
\midrule
\(GF(4)\) base isomorphism & fixed algebraic convention\\
Sparse allele residual \(\delta\) & falsifiable, allele-specific\\
Syndrome \(Hx=B\delta\) & constrained algebraic program\\
Toeplitz background operators & collision remediation for \(x\)\\
Unique molecule inverse from \(x\) & \textbf{not claimed}\\
Docking / affinity / steric clearance & \textbf{not claimed}\\
Clinical efficacy or safety & \textbf{not claimed}
(\texttt{clinical\_grade=false})\\
\bottomrule
\end{tabular}
\end{center}

Algebraic collision pathology is addressed at the tensor level; transition to
3D pose generation remains a separate geometric observer lane and is outside
the scope of this document.

\subsection*{Appendix B disclaimer}
This document is a computational-chemistry method ledger only. It encodes
finite-field allele isolation and background-conditioned syndrome solving. It
does not claim therapeutic effect, manufacturing readiness, or empirical
binding.



\end{document}
